\section{From the constructed law to the matrix-exponential semigroup}
\label{sec:semigroup}

\subsection{The terminal marginal is a genuine pushforward}

For arithmetic convenience, the development first defines the sector mass that ends at $y$ and sums it over $n$ to obtain a transition mass $p_T(x,y)$. A separate bridge theorem proves that this sum is the terminal marginal of the actual path law:
\begin{equation}
\label{eq:terminal-push}
  (\Pp_{T,x}\circ X_T^{-1})(\{y\})=p_T(x,y),
\end{equation}
where $X_T$ is the terminal-state coordinate. This step matters in a proof assistant: without it, a calculation about $p_T$ would not automatically be a statement about \lean{pathLawFrom}.

\subsection{Renewal equation and uniqueness}

The path construction satisfies a first-jump renewal equation. The Lean proof is formulated in a residual fraction $\rho\in[0,1]$ of a fixed horizon. Peeling the first free simplex coordinate leaves the remaining coordinates in a simplex at fraction $\rho-u_0$ without rescaling the chart. After summing over the first successor state, the transition mass satisfies the same entrywise renewal equation as the matrix exponential.

On the analytic side, the development proves that
\[
  \rho\longmapsto \exp((T\rho)Q)(x,y)
\]
is a continuous solution of that renewal equation and that bounded continuous solutions are unique. Matching the path-side solution to the matrix-exponential solution gives
\begin{equation}
\label{eq:terminal-exp}
  (\Pp_{T,x}\circ X_T^{-1})_{\R}(\{y\})
  = \exp(TQ)(x,y),
\end{equation}
where the subscript indicates conversion from an extended-nonnegative-real measure to its real-valued mass: the left-hand side is a finite singleton mass, read as a real number.

In outline, the proof of~\eqref{eq:terminal-exp} runs as follows. Conditioning on whether the path jumps within the residual fraction expresses $p_{T\rho}(x,y)$ as a no-jump survival term plus an integral over the first holding time and first successor of the same transition mass at a smaller fraction; this is the path-side renewal equation. The transition mass is bounded by one and continuous in $\rho$, and the matrix-exponential entry solves the same equation. For uniqueness, the difference of two bounded continuous solutions satisfies a homogeneous integral inequality; iterating it $n$ times bounds the difference by $C(Rt)^n/n!$, which tends to zero, so the solutions coincide. Combining the identification with the bridge~\eqref{eq:terminal-push} transfers the equality to the constructed law.

The terminal laws are packaged as a Mathlib Markov kernel
\[
  K_T(x,\cdot)=\Pp_{T,x}\circ X_T^{-1}.
\]
Equation~\eqref{eq:terminal-exp} turns the matrix-exponential semigroup identity into Chapman--Kolmogorov:
\[
  K_{S+T}=K_T\circ K_S.
\]
Here the rightmost kernel acts first, following Mathlib's kernel-composition convention. The repository also proves a stronger path-level identity. Cutting a constructed path at an intermediate horizon and concatenating an independent continuation started at the cut state recovers \lean{pathLawFrom (S+T) x}. Taking the terminal pushforward of this path-level theorem yields the kernel semigroup law as a shadow.

\subsection{Finite-dimensional distributions}

Let
\[
  0=t_0\le t_1\le\cdots\le t_n\le T
\]
be a finite monotone family of observation times. The measurable sampling map reads the right-continuous trajectory at these times. The formalization proves
\begin{equation}
\label{eq:fdd}
  \Pp_{T,x}\circ (X_{t_0},\ldots,X_{t_n})^{-1}
  = \delta_x K_{t_1-t_0}\cdots K_{t_n-t_{n-1}},
\end{equation}
where the right-hand side is Mathlib's chronological finite path measure. For a state sequence $(x_0,\ldots,x_n)$, the atom formula is
\begin{equation}
\label{eq:fdd-atom}
  \mathbf{1}_{\{x_0=x\}}
  \prod_{i=0}^{n-1}
  \exp((t_{i+1}-t_i)Q)(x_i,x_{i+1}).
\end{equation}
Thus the semigroup is not merely an API defined beside the path construction: it is identified with the actual terminal and finite-dimensional laws of that construction.

A branching three-state Y-shaped chain instantiates the general theorem. This example is useful because its path law cannot be reduced to the parity-of-jump-count argument available for a two-state alternating chain.

